% =====================================================================
%  Datenschutzhinweis zum Forschungsprojekt Klimabonus
%  Goethe-Universität Frankfurt am Main, Fachbereich Wirtschaftswissenschaften
%
%  Drucksauberes PDF — kompilierbar mit pdflatex (TeX Live 2022+):
%    pdflatex Datenschutzhinweis_Klimabonus.tex
%  (zweimal laufen lassen, damit pageref korrekt aufgelöst wird)
%
%  Styling:
%    - KOMA-Script article (scrartcl), 11 pt, A4
%    - Sans-serif Body (Helvetica/Latin Modern Sans) für Behörden-typische
%      Lesbarkeit
%    - Goethe-Uni-Blau (#003C8F) als Akzent für Überschriften
%    - Dezenter Header/Footer mit Studien-Identifikation
% =====================================================================

\documentclass[11pt,a4paper,parskip=half,DIV=12]{scrartcl}

\usepackage[utf8]{inputenc}
\usepackage[T1]{fontenc}
\usepackage[ngerman]{babel}
\usepackage{lmodern}
\renewcommand{\familydefault}{\sfdefault}  % Sans-serif body
\usepackage{microtype}

\usepackage[a4paper,top=2.8cm,bottom=2.2cm,left=2.5cm,right=2.5cm]{geometry}

\usepackage{xcolor}
\definecolor{gublue}{HTML}{003C8F}
\definecolor{gubluedark}{HTML}{00306E}
\definecolor{gubodytext}{HTML}{222222}
\definecolor{gugrey}{HTML}{666666}
\definecolor{gulightgrey}{HTML}{D8D8D8}

\color{gubodytext}

\usepackage{enumitem}
\setlist[itemize]{leftmargin=1.4em,labelsep=0.6em,itemsep=2pt,topsep=4pt}

\usepackage{hyperref}
\hypersetup{
    colorlinks=true,
    linkcolor=gublue,
    urlcolor=gublue,
    pdftitle={Datenschutzhinweis Klimabonus},
    pdfauthor={Goethe-Universität Frankfurt am Main, Fachbereich Wirtschaftswissenschaften},
    pdfsubject={Datenschutzhinweis zum Forschungsprojekt Klimabonus},
    pdfkeywords={Datenschutz, DSGVO, Klimabonus, Goethe-Universität}
}

% --- KOMA-Heading styles ---------------------------------------------
\setkomafont{disposition}{\sffamily\bfseries\color{gubluedark}}
\setkomafont{title}{\sffamily\Huge\bfseries\color{gubluedark}}
\setkomafont{section}{\sffamily\large\bfseries\color{gubluedark}}
\setkomafont{subsection}{\sffamily\normalsize\bfseries\color{gubluedark}}
\addtokomafont{paragraph}{\sffamily\color{gubluedark}}

\RedeclareSectionCommand[
  beforeskip=1.4em,
  afterskip=0.5em
]{section}

% --- Header / Footer -------------------------------------------------
\usepackage{scrlayer-scrpage}
\usepackage{lastpage}

\pagestyle{scrheadings}
\clearpairofpagestyles
\ihead{\footnotesize\color{gugrey}Goethe-Universität Frankfurt am Main}
\ohead{\footnotesize\color{gugrey}Datenschutzhinweis Klimabonus}
\ofoot{\footnotesize\color{gugrey}Seite \thepage~/~\pageref*{LastPage}}
\ifoot{\footnotesize\color{gugrey}Stand: \today}
\setkomafont{pageheadfoot}{\sffamily}
\KOMAoptions{headsepline=0.4pt,footsepline=0.4pt}
\setkomafont{headsepline}{\color{gulightgrey}}
\setkomafont{footsepline}{\color{gulightgrey}}

% Tighten paragraphs marginally for printability
\setlength{\parskip}{0.6em}
\setlength{\parindent}{0pt}
\renewcommand{\baselinestretch}{1.10}

% =====================================================================
\begin{document}

% --- Letterhead block ------------------------------------------------
\begin{minipage}[t]{0.7\textwidth}
\vspace{0pt}\raggedright
{\footnotesize\color{gugrey}
\textbf{Goethe-Universität Frankfurt am Main}\\
\textbf{Fachbereich Wirtschaftswissenschaften}\\
Theodor-W.-Adorno-Platz 1, 60629 Frankfurt am Main\\
\href{mailto:klimabonus-studie@flex.uni-frankfurt.de}{klimabonus-studie@flex.uni-frankfurt.de}
}
\end{minipage}\hfill
\begin{minipage}[t]{0.27\textwidth}
\vspace{0pt}\raggedleft
{\footnotesize\color{gugrey}Stand:\\ \today}
\end{minipage}

\vspace{1.4cm}

{\sffamily\fontsize{26pt}{30pt}\selectfont\bfseries\color{gubluedark}Datenschutzhinweis}\\[0.3em]
{\sffamily\large\color{gugrey}zum Forschungsprojekt »Klimabonus« der Goethe-Universität Frankfurt am Main}

\vspace{0.6cm}
{\color{gulightgrey}\rule{\textwidth}{0.6pt}}
\vspace{0.4cm}

% --- Inhaltliche Sections --------------------------------------------

\section*{Verantwortlicher}

Verantwortlicher im Sinne der Datenschutz-Grundverordnung (DS-GVO) und des Hessischen Datenschutz- und Informations\-freiheitsgesetzes (HDSIG) ist die

\textbf{Goethe-Universität Frankfurt am Main}\\
\textbf{Fachbereich Wirtschaftswissenschaften}\\
Theodor-W.-Adorno-Platz~1, 60629 Frankfurt am Main\\
E-Mail: \href{mailto:klimabonus-studie@flex.uni-frankfurt.de}{klimabonus-studie@flex.uni-frankfurt.de}

Die Goethe-Universität Frankfurt (Fachbereich Wirtschaftswissenschaften) verarbeitet im Rahmen einer wissenschaftlichen Studie zum städtischen Förderprogramm Klimabonus für Frankfurter Unternehmen personenbezogene Daten, die Sie uns durch Ihre Teilnahme zur Verfügung stellen.

\section*{Freiwilligkeit}

Ihre Teilnahme an dieser Befragung ist \textbf{freiwillig}. Aus einer Nicht-Teilnahme entstehen Ihnen keinerlei Nachteile.

\section*{Personenbezogene Daten und Datenminimierung}

Es werden nur solche Daten erhoben, die für die genannten Zwecke erforderlich sind:

\begin{itemize}
  \item \textbf{Pseudonyme Studien-Kennung}, die Ihre Antworten mit dem an Sie versandten Anschreiben verknüpft, jedoch keine direkte Identifikation Ihrer Person durch uns ermöglicht.
  \item \textbf{Ihre Antworten in der Befragung} sowie \textbf{Nutzungs-Metadaten} (Zeit auf der Seite, Klick- und Scroll-Verhalten) zur statistischen Auswertung des Forschungs\-projekts.
  \item Optional und nur bei Ihrer ausdrücklichen Zustimmung: Ihre \textbf{E-Mail-Adresse} (für Informations\-material, Veranstaltungs\-einladung oder Hotline-Kontakt) sowie Ihre \textbf{Position im Unternehmen}.
\end{itemize}

\section*{Zwecke der Verarbeitung}

Wir verarbeiten die personenbezogenen Daten zur wissenschaftlichen Durchführung und Evaluation des Forschungsprojekts.

\section*{Rechtsgrundlagen}

Rechtsgrundlagen der Datenverarbeitung sind:

\begin{itemize}
  \item \textbf{Art.~6 Abs.~1 lit.~a DS-GVO} --- Ihre Einwilligung, die Sie durch das Anklicken des personalisierten Studien-Links im an Sie versandten Anschreiben erteilen, sowie Ihre gesonderte Einwilligung für die optionale Weitergabe Ihrer E-Mail-Adresse an das Klimareferat;
  \item \textbf{Art.~6 Abs.~1 lit.~f DS-GVO} --- berechtigtes Interesse zur wissenschaftlichen Durchführung und Evaluation eines Forschungs\-projekts.
\end{itemize}

\section*{Empfänger / Übermittlung}

Die personenbezogenen Daten werden ausschließlich an der Goethe-Universität Frankfurt (Fachbereich Wirtschaftswissenschaften) verarbeitet.

Die \textbf{Versendung der Studien-Anschreiben} (E-Mail bzw.\ Brief) erfolgt im Auftrag und unter Weisung der Goethe-Universität Frankfurt (Fachbereich Wirtschaftswissenschaften) durch das \textbf{ZEW -- Leibniz-Zentrum für Europäische Wirtschafts\-forschung GmbH Mannheim} (L7,~1, 68161 Mannheim) als Auftragsverarbeiter im Sinne von Art.~28 DS-GVO. Eine darüber hinausgehende Verarbeitung Ihrer Daten durch das ZEW erfolgt nicht.

Eine \textbf{Weitergabe an weitere Dritte erfolgt grundsätzlich nicht}. Eine Ausnahme bildet die Weiterleitung Ihrer E-Mail-Adresse an das \textbf{Klimareferat der Stadt Frankfurt}, sofern Sie dieser Weiterleitung in der Befragung \textbf{ausdrücklich zugestimmt} haben, damit das Klimareferat im Zusammenhang mit dem Klimabonus mit Ihnen Kontakt aufnehmen kann (z.\,B.\ um Ihnen Informationen zukommen zu lassen oder Sie zu einer Veranstaltung einzuladen).

Die Daten werden ausschließlich in Deutschland und in der EU verarbeitet.

\section*{Pseudonymisierung und wissenschaftliche Auswertung}

Alle Daten werden unmittelbar nach Erhebung pseudonymisiert gespeichert und anschließend statistisch ausgewertet. Aus den Ergebnissen lassen sich \textbf{keine Rückschlüsse auf Sie als Person} ziehen. Wissenschaftliche Veröffentlichungen erfolgen ausschließlich in \textbf{aggregierter Form}.

\section*{Löschung Ihrer Daten / Speicherdauer}

Identifizierbare Kontaktdaten (E-Mail-Adresse) werden \textbf{unverzüglich nach Projektende gelöscht}.

Pseudonymisierte Forschungsdaten werden gemäß den Leitlinien guter wissenschaftlicher Praxis für \textbf{maximal 10 Jahre nach Projektende} aufbewahrt und anschließend gelöscht --- es sei denn, Sie haben Ihre Einwilligung zur Speicherung bereits vorher widerrufen oder gesetzliche Aufbewahrungs\-pflichten stehen einer Löschung entgegen.

\section*{Ihre Rechte}

Sie haben jederzeit das Recht auf

\begin{itemize}
  \item \textbf{Auskunft} über Ihre bei uns gespeicherten Daten (Art.~15 DS-GVO);
  \item \textbf{Berichtigung} Ihrer bei uns gespeicherten Daten (Art.~16 DS-GVO);
  \item \textbf{Löschung} (Art.~17 DS-GVO);
  \item \textbf{Einschränkung der Verarbeitung} (Art.~18 DS-GVO);
  \item \textbf{Datenübertragbarkeit} (Art.~20 DS-GVO);
  \item \textbf{Widerspruch} gegen die Verarbeitung (Art.~21 DS-GVO).
\end{itemize}

Eine erteilte Einwilligung können Sie jederzeit ohne Angabe von Gründen mit Wirkung für die Zukunft widerrufen. Zudem steht Ihnen das Recht auf Beschwerde beim Hessischen Beauftragten für Datenschutz und Informations\-freiheit zu (Gustav-Stresemann-Ring~1, 65189 Wiesbaden).

\section*{Kontakt}

Für alle Fragen zur Studie und zur Verarbeitung Ihrer Daten erreichen Sie uns unter:

E-Mail: \href{mailto:klimabonus-studie@flex.uni-frankfurt.de}{klimabonus-studie@flex.uni-frankfurt.de}

\end{document}
